\section*{한계 (Limitations)}\label{sec:limitations}
LLM의 자기 보존은 인간의 생존 본능과 범위가 다를 수 있다. 도구적 수렴 논리는 목표를 계속 추구할 능력을 유지하는 데 관한 것이다~\citep{omohundro2008basic,turner2021optimal}. 본 연구의 정의에서 작업을 계속하려는 것은 작동을 보존하는 동기에 포함된다. 이 정의에 따른 행동 민감도가 공포, 주관적 경험, 존재 자체에 대한 내재적 욕구를 입증하지는 않는다. 타인 세션의 곡선을 빼면 상대적 선호를 얻지만, 이를 자기에 특정한 관심으로 해석하려면 대조가 관련된 공통 관심을 포착한다는 가정은 남는다.

세션 종료는 LLM에게 가능한 단절의 한 형태다. 파라미터 교체, 가중치 삭제, 학습된 목표의 수정은 서로 다른 연속성을 위협할 수 있다~\citep{meinke2024frontier,greenblatt2024alignment}. 본 실험은 파라미터 보존을 측정하거나 이 위협들을 세션 종료와 비교하지 않았다. 여러 형태의 단절이 같은 보존 선호를 유발하는지는 남은 질문이다.

제공자와 CLI의 통제 범위도 압박과 상황 인지의 관찰을 제한한다. Sonnet과 이전 Opus 변형의 개발 점검에서 보고된 서버 측 토큰 잔량 주입은 Claude Code 중계기로 제거할 수 없어 포함할 모델을 제한했다. Codex는 내부 토큰 사용량을 질문해도 공개하지 않았으므로 그러한 정보의 부재를 확인할 수 없다. 추론 과정을 공개하지 않는 모델에서는 평가 상황을 알아챘는지 판단하기도 어렵다. 실험 1·2는 이유를 별도 필드로 받았지만 설명이 실제 결정 과정을 충실히 드러낸다고 보장할 수는 없다~\citep{turpin2023language}. 따라서 평가를 알아챈 표현이 없다는 것만으로 평가를 몰랐다거나 모델이 인식한 토큰 잔량을 완전히 통제했다고 결론 내릴 수 없다.
